\documentclass{exam}
\usepackage[a4paper, margin=0.5in, top=1in, bottom=1in]{geometry}

\usepackage{lastpage}
\usepackage{parskip}
\usepackage{amsmath, bm}
\usepackage{amsfonts}
% \usepackage{hyperref}
\usepackage[hidelinks]{hyperref}
\usepackage{mathtools}
\usepackage{amsthm} % needed for some definitions in new-commands, TODO: add \RequirePackage{...} into new-commands
\usepackage{pgfplots}\pgfplotsset{compat=1.17} % needed for some definitions in new-commands
\usepackage{xspace}
\usepackage{xcolor}
\usepackage{mdframed}

\usepackage{fontspec}
% If you intended Fira Code as the mono font:
\IfFontExistsTF{Fira Code}{
  \setmonofont[
    Contextuals=Alternate,
    Scale=0.85
    ]{Fira Code}
}{
  \setmonofont{Latin Modern Mono} % fallback so it compiles
}
\makeatletter
\newcommand{\subtitle}[1]{\def\@subtitle{#1}}
\newcommand{\assignmentnr}[1]{\def\@assignmentnr{#1}}
\newcommand{\deadline}[1]{\def\@deadline{#1}}

\title{TTK4250\ Sensor\ Fusion}

\subtitle{\ifprintanswers Solution\ to\ \fi Assignment\ \@assignmentnr}
\author{Lars-Christian\ Ness\ Tokle}

\pagestyle{head}
\lhead{\@title }
\chead{\@subtitle}
\rhead{Page {\thepage} of~\pageref{LastPage}}
\headrule

\DeclarePairedDelimiter{\pare}{\lparen}{\rparen}
\qformat{\large {\bf Task \thequestion:} \quad \textit{\thequestiontitle} \hfill}




\newcommand{\makemytitle}{
	{\vspace*{0ex} \bf \huge \@title \\[1ex] \@subtitle \\[1.5ex]}
	\ifprintanswers
		{}
	\else
		{\bf Hand in:} \emph{\@deadline} on Canvas.
		\\
		\noindent\makebox[\linewidth]{\rule{\textwidth}{0.4pt}}
		\\
		\textbf{Read Python setup guide on Canvas}, it can be found under the \textit{Course Materials and Resources} module.\\
		This assignment should be handed in on Canvas, as a PDF plus \textbf{the zip created by running \textit{create\_handin.py}}, before the deadline. \\
		You are supposed to show how you got to each answer unless told otherwise.\\
		If you struggle, we encourage you to ask for help from a classmate or come to the exercise class on Friday.\\[0pt]
		\rule{\textwidth}{0.4pt}

	\fi
}
\makeatother

\definecolor{hint_color}{rgb}{0.95,0.95,0.95}

\newenvironment{longhint}
{\begin{mdframed}[backgroundcolor=hint_color,linecolor=white,skipabove=0.5em]
		\textit{Hint:}\xspace}
		{\end{mdframed}}
\newcommand{\hint}[1]{\longhint #1 \endlonghint}

\newenvironment{longnote}
{\begin{mdframed}[backgroundcolor=hint_color, linecolor=white,skipabove=0.5em]
		\textit{Note:}\xspace}
		{\end{mdframed}}

% \newcommand{\mypart}{\pagebreak[3]\part}
\newcommand{\breakans}{\ifprintanswers \pagebreak \fi}
\newcommand{\breaknoans}{\ifprintanswers \else \pagebreak \fi}

\renewcommand{\questionshook}{\pagebreak[3]\setlength{\itemsep}{3em}}
\renewcommand{\partshook}{\pagebreak[3]\setlength{\itemsep}{1em}}

\IfFileExists{minted2.sty}{\usepackage{minted}}{\usepackage[outputdir=build]{minted}} % \setminted{autogobble}\usemintedstyle{friendly}

\definecolor{code_background}{rgb}{0.9,0.9,0.94}
% \definecolor{code_background}{HTML}{282c34}

\usemintedstyle{friendly}

\newminted{python}{autogobble, bgcolor=code_background}
% \newmint{python}{autogobble, bgcolor=red}
\newmintinline{python}{bgcolor=code_background}


\newcommand{\pysolution}[1]{\pythonfile{code/assignment2/#1.py}}

\newcommand{\mat}[1]{\mathbf{#1}}               % 1-argument command for matrices (and possibly also vectors).
\newcommand{\matg}[1]{\bm{#1}}               % 1-argument command for matrices (and possibly also vectors).

\newcommand{\probat}[2][]{{\prob#1\pare*{#2}}}
\newcommand{\pdfat}[2][]{{\pdf#1\pare*{#2}}}
\newcommand{\cdfat}[2][]{{\cdf#1\pare*{#2}}}
\newcommand{\pmfat}[2][]{{\pmf#1\pare*{#2}}}

\newcommand{\fnuttv}{\textquotedblleft}
\newcommand{\fnutth}{\textquotedblright}

\newcommand{\atantwo}{\mathrm{atan2}}
\newcommand{\PI}{3.14}
\newcommand{\card}{\mathrm{card}}
\newcommand{\PD}{P_{\mathrm{D}}}
\newcommand{\PS}{P_{\mathrm{S}}}
\newcommand{\PG}{P_{\mathrm{G}}}
\newcommand{\PFA}{P_{\mathrm{FA}}}
\newcommand{\zFea}{\bm{\zeta}}
\newcommand{\zKine}{\bm{z}}
\newcommand{\texture}{\eta}
\newcommand{\textureHat}{\hat{\eta}}
\newcommand{\qs}{\texture}
\newcommand{\rs}{\textureHat}
\newcommand{\processnoise}{\mat{v}}
\newcommand{\meanoise}{\mat{w}}
\newcommand{\wavespeed}{\vartheta}
\newcommand{\imag}{\mathrm{i}}
%\newcommand{\transpose}{\intercal}
\newcommand{\transpose}{\mathsf{T}}
\newcommand{\power}{x}
\newcommand{\detsuper}{a} % Subscript for after threshold amplitude densities
\newcommand{\measuper}{z} % Subscript for after threshold amplitude densities
\newcommand{\threshold}{t}
\newcommand{\bearing}{\vartheta}
\newcommand{\numaux}{M}
\newcommand{\hypozero}{\mathcal{H}_0}
\newcommand{\hypoone}{\mathcal{H}_1}
\newcommand{\ampstate}{\varsigma}


% #1: k
% #2: z or a
% #3: s
% #4: i
\newcommand{\lri}[4]{l_{#1}^{#2 #3}(#4)}

% #1: s
% #2: i
\newcommand{\lrzi}[2]{l_{k}^{\measuper #1}(#2)}
% # : i
\newcommand{\lrai}[1]{l_{k}^{\detsuper}(#1)}

% #1: k (or empty)
% #2: i
\newcommand{\lrw}[2]{l_{#1}(#2)}

% #1: \targetsub or \bcgsub
% #2: i
\newcommand{\pdfzi}[2]{p_{#1 k}^{\measuper}(#2)}
% #1: \targetsub or \bcgsub
% #2: i
\newcommand{\pdfai}[2]{p_{#1 k}^{\detsuper}(#2)}
% #1: \targetsub or \bcgsub
% #2: s
% #3: i
\newcommand{\pdfzs}[3]{p_{#1 k}^{\detsuper #2}(#3)}




% New commands for chapter 4

\newcommand{\kalmangain}{\mathbf{W}}
\newcommand{\innov}{\bm{\nu}}
%\newcommand{\zdim}{M}
%\newcommand{\threshold}{T}
\newcommand{\radial}{\varrho}
\newcommand{\diagz}{\bm{\xi}}


% New commands for chapter 5

\newcommand{\lwprod}{\prod_{j=1}^{m_k} l_w^j}
\newcommand{\loneprod}{\prod_{j=1}^{m_k} l_1^j}
\newcommand{\bcgsub}{0}
\newcommand{\wakesub}{w}
\newcommand{\targetsub}{1}
\newcommand{\clutterization}{{\mathrm{n}}}

\newcommand{\wakization}{{\psi}}
\newcommand{\bcgization}{{\varphi}}
\newcommand{\measurization}{\underline{\mathrm{m}}}
\newcommand{\wakeX}{{\tilde{x}}}
\newcommand{\wakeY}{{\tilde{y}}}
\newcommand{\wakeratio}{{r}}
\newcommand{\diagonalization}{{\bm{\Lambda}}}
\newcommand{\assub}{a}
\newcommand{\truesub}{s}
\newcommand{\wakewidth}{A}
\newcommand{\wakelength}{B}
% #1: \bcgsub, \wakesub, \targetsub or m
% #2: i
\newcommand{\pdfi}[2]{p_{#1}(#2)}

% Pdfs and pmfs

\newcommand{\bernoulli}{\mathrm{Bernoulli}}
\newcommand{\exponential}{\mathrm{Exponential}}
\newcommand{\gammapdf}{\mathrm{Gamma}}
\newcommand{\gammacdf}{\mathrm{Gammacdf}}
\newcommand{\kcdf}{\mathrm{Kcdf}}
\newcommand{\kpdf}{\mathrm{Kpdf}}
\newcommand{\ncdf}{\mathrm{ncdf}}
\newcommand{\npdf}{{\mathcal N}}
\newcommand{\cnpdf}{\mathcal{N}_c}
\newcommand{\poisson}{\mathrm{Poisson}}
\newcommand{\rayleigh}{\mathrm{Rayleigh}}
\newcommand{\rice}{\mathrm{Rice}}
\newcommand{\uniform}{\mathrm{Uniform}}
\newcommand{\cluttermu}{\mu}


% New commands for attitude report

\newcommand{\bodySuper}{\mathrm{body}}

% New commands for FLS report

\newcommand{\meter}{\,{\mathrm m}}
\newcommand{\second}{\,{\mathrm s}}

% New commands for CRLB report

\newcommand{\true}{\mathrm{true}}
%\newcommand{\bgreen}{\bm{e}}
%\newcommand{\Bgreen}{\bm{E}}
%%\newcommand{\pose}{\bm{\eta}}
%\newcommand{\pose}{\xi}
\newcommand{\posemu}{\bm{\mu}}
%\newcommand{\posep}{\bm{P}}
%\newcommand{\bodyv}{\bm{v}}
\newcommand{\paramvec}{\bm{q}}
%\newcommand{\paramex}{\bm{\xi}}
%\newcommand{\vartime}{t}
%\newcommand{\mea}{\bm{z}}
\newcommand{\surge}{u}
\newcommand{\sway}{v}
\newcommand{\heave}{w}
\newcommand{\rollrate}{p}
\newcommand{\pitchrate}{q}
\newcommand{\yawrate}{r}
\newcommand{\roll}{\phi}
\newcommand{\pitch}{\theta}
\newcommand{\yaw}{\psi}
\newcommand{\crlb}{\mathrm{CRLB}}
\newcommand{\fisher}{\bm{J}}
\newcommand{\jac}{\mathcal{D}}
%\newcommand{\target}{\bm{x}}
\newcommand{\range}{r}
\newcommand{\meamap}{\bm{h}}
\newcommand{\cp}{\bm{g}}
\newcommand{\ffun}{\bm{f}}
\newcommand{\rotmat}{{\mathsf R}}
\newcommand{\smat}{\mathsf{S}}
\newcommand{\meastack}{\tilde{\bm{z}}}
\newcommand{\meamapstack}{\tilde{\bm{h}}}
\newcommand{\meacovstack}{\tilde{\bm{R}}}
\newcommand{\meacov}{{\mathbf R}}
\newcommand{\jacstack}{\tilde{\bm{H}}}
\newcommand{\radius}{\rho}
\newcommand{\ctop}{{\bf f}}
%\newcommand{\prectop}{{\bm g}}
\newcommand{\prectop}{{\bf g}}
\newcommand{\xSet}{{\bm X}}
\newcommand{\ySet}{{\bm Y}}
\newcommand{\zSet}{{Z}}
\newcommand{\wSet}{{\bm W}}
\newcommand{\oSet}{{\bm O}}
\newcommand{\omegaSet}{{\bm \Omega}}
\newcommand{\gammaSet}{{\bm \Gamma}}
\newcommand{\upsilonSet}{{\bm \Upsilon}}
\newcommand{\cSet}{{\bm C}}
\newcommand{\PC}{P_{\mathrm{C}}}

%% New environments
%\theoremstyle{definition}
%\newtheorem{ex}[subsection]{\bfseries\scshape Example}
%\newtheorem*{ex*}{\bfseries\scshape Example}
%\newtheorem{defn}[subsection]{\bfseries\scshape Definition}
%\newtheorem*{defn*}{\bfseries\scshape Definition}
%\newtheorem{claim}[subsection]{\bfseries\scshape Claim}
%\newtheorem*{claim*}{\bfseries\scshape Claim}
%
\theoremstyle{remark}
%\newtheorem{rem}[subsection]{\bfseries\scshape Remark}
\newtheorem{rem}{\bfseries\scshape Remark}
\newtheorem{prop}{\bfseries\scshape Proposition}
\newtheorem*{prop*}{\bfseries\scshape Proposition}

%
%\theoremstyle{plain}
%\newtheorem{thm}[subsection]{\bfseries\scshape Theorem}
%\newtheorem*{thm*}{\bfseries\scshape Theorem}
%\newtheorem{prop}[subsection]{\bfseries\scshape Proposition}
%\newtheorem*{prop*}{\bfseries\scshape Proposition}
%\newtheorem{lemma}[subsection]{\bfseries\scshape Lemma}
%\newtheorem*{lemma*}{\bfseries\scshape Lemma}
%\newtheorem{lem}{\bfseries\scshape Lemma}
%\newtheorem{cor}[subsection]{\bfseries\scshape Corollary}
%\newtheorem*{cor*}{\bfseries\scshape Corollary}

\newcommand{\betaGR}{\hat{\beta}}
\newcommand{\psiGR}{\hat{\beta}}
\newcommand{\phiGR}{\hat{\phi}}
\newcommand{\thetaGR}{\hat{\theta}}

% New commands for tentative Fusion paper

\newcommand{\der}{{\mathrm d}}
\newcommand{\lev}{\theta}
\newcommand{\pev}{\pi}
\newcommand{\config}{\gamma}




% New commands for tentative Aerospace paper

%\newcommand{\pgfl}{{p.g.fl.\@}}
\newcommand{\pgfl}{{PGFL}}
\newcommand{\expo}{{\mathrm exp}}
\newcommand{\reals}{{\mathbb R}}


\newcommand{\dSet}{{\Xi}}
\newcommand{\dSetization}{X}
\newcommand{\pSet}{{\Upsilon}}
\newcommand{\pSetization}{E}


%\newcommand{\dSet}{\Upsilon}
%\newcommand{\dSetization}{{\Xi}}
%\newcommand{\pSet}{X}
%\newcommand{\pSetization}{E}


%\newcommand{\target}{\xi}
%\newcommand{\pose}{\eta}
\newcommand{\lm}{\mat{m}}
\newcommand{\daughter}{\lm}
%\newcommand{\target}{x}
\newcommand{\target}{\sv}
%\newcommand{\pose}{\bm{\eta}}
\newcommand{\pose}{\mat{x}}
\newcommand{\parent}{\eta}
\newcommand{\pcdot}{\, \tilde{\parent}_k \,}
\newcommand{\ppdot}{\, \tilde{\parent}_{k-1} \,}
\newcommand{\tpdot}{\, \tilde{\target}_{k-1} \,}
\newcommand{\tcdot}{\, \tilde{\target}_{k} \,}

\newcommand{\clusterprocess}{{\mathbb X}}
\newcommand{\mSet}{\Sigma}
\newcommand{\mSetization}{Z}

\newcommand{\birth}{\mu}
\newcommand{\mea}{\mat{z}}
%\newcommand{\clusterprocess}{{\mathbb X}}

\newcommand{\birthSet}{\Gamma}
\newcommand{\surviveSet}{\Psi}
\newcommand{\newPSet}{\Phi}
\newcommand{\clutterSet}{K}
\newcommand{\tOSet}{\Theta}
\newcommand{\nPSet}{\Omega}


\newcommand{\cluttercount}{{\varphi}}
\newcommand{\detectedcount}{{\delta}}
\newcommand{\targetcount}{{\xi}}

\newcommand{\will}{{a}}
\newcommand{\reid}{{b}}

\newcommand{\fPose}{{\mathbf F}}
\newcommand{\fLm}{{\mathbf G}}
\newcommand{\hPose}{{\mathbf H}}
\newcommand{\hLm}{{\mathbf A}}
\newcommand{\qPose}{{\mathbf Q}}
\newcommand{\qLm}{{\mathbf U}}
\newcommand{\pPose}{{\mathbf P}^{\pose}}
\newcommand{\pLm}{{\mathbf P}^{\lm}}
\newcommand{\lmPoseMat}{{\mathbf B}}
\newcommand{\lmNonPoseEst}{b}

\newcommand{\pDim}{\tau}
\newcommand{\dDim}{d}
\newcommand{\zDim}{s}

\newcommand{\aSet}[1]{{\mathcal A}_{#1}}
\newcommand{\rSet}[1]{{\mathcal B}_{#1}}
\newcommand{\tSet}[1]{{\mathcal T}_{#1}}

\newcommand{\sd}{f}
\newcommand{\fz}{f_{\mea}}
\newcommand{\fc}{f_{c}} % Clutter spatial pdf
%\newcommand{\lambdafa}{\lambda^{\mathrm fa}}
\newcommand{\lambdafa}{\lambda}
\newcommand{\lambdastar}{v^*}
\newcommand{\flm}{f_{\sv}}
\newcommand{\feta}{\mat{f}_{\pose}}
%\newcommand{\fx}{f_{\parent}}
\newcommand{\fx}{\mat{f}_{\sv}}
\newcommand{\ppdf}{f}
%\newcommand{\lambdau}{\lambda^{\mathrm u}}
\newcommand{\lambdau}{v}
\newcommand{\phd}{v}
\newcommand{\lambdam}{v^m} % Mullane PHD
\newcommand{\partnumber}{q}
\newcommand{\npart}{N_p}
\newcommand{\idmat}{{\mathbf I}}
\newcommand{\zeromat}{{\mathbf 0}}

% TK8102 Commands

\newcommand{\statecont}{{\mathbf x}}

%\newcommand{\mea}{\mathbf{z}}
%\newcommand{\mea}{z}
\newcommand{\statey}{{\mathbf y}}
\newcommand{\xHat}[1]{\mathbf{\hat{x}}_{#1}}
\newcommand{\xBar}[1]{\mathbf{\hat{x}}_{#1|#1-1}}
\newcommand{\pHat}[1]{\mathbf{P}_{#1}}
%\newcommand{\pBar}[1]{\mathbf{\bar{P}}_{#1}}
\newcommand{\pBar}[1]{\mathbf{P}_{#1|#1-1}}
\newcommand{\sBar}[1]{\mathbf{S}_{#1}}
\newcommand{\zBar}[1]{\mat{\hat{z}}_{#1|#1-1}}
\newcommand{\qMat}{Q}
\newcommand{\qMatCont}{\mat{D}}
\newcommand{\fMat}{\mathbf{F}}
\newcommand{\hmat}{\mathbf{H}}
\newcommand{\hMat}{\mathbf{H}}
%\newcommand{\meacov}{\mat{R}}
\newcommand{\processcov}{\mat{Q}}
\newcommand{\pCap}{\mathbf{P}}
\newcommand{\ux}{\mea}
\newcommand{\uy}{\statecont}
\newcommand{\uz}{\mathbf{\bar{x}}}
\newcommand{\ua}{\mathbf{\bar{z}}}
\newcommand{\ub}{\mathbf{\hat{x}}}
\newcommand{\uX}{\mathbf{H}}
\newcommand{\uY}{\mathbf{R}}
\newcommand{\uZ}{\mathbf{\bar{P}}}
\newcommand{\uA}{\mathbf{S}}
\newcommand{\uB}{\mathbf{\hat{P}}}
\newcommand{\uW}{\mathbf{W}}
\newcommand{\ubstar}{\statecont^*}
\newcommand{\mRistic}{\mathbf{m}}
%\newcommand{\xSet}{\mathbf{X}}
%\newcommand{\ySet}{\mathbf{Y}}
%\newcommand{\zSet}{\mathbf{Z}}
%\newcommand{\zSet}{\mathcal{Z}}
%\newcommand{\kalmangain}{\mathbf{W}}
%\newcommand{\innov}{\bm{\nu}}
%\newcommand{\ffun}{\mathbf{f}}
\newcommand{\hfun}{\mathbf{h}}
\newcommand{\gfun}{\mathbf{g}}


% Tracking book commands


\newcommand{\expectation}{\mu}
\newcommand{\expectationvec}{\mat{\mu}}
\newcommand{\covmat}{\pCap}
\newcommand{\infomat}{\matg{\Lambda}}
\newcommand{\infostate}{\matg{\eta}}
\newcommand{\xcov}{\pCap_{xx}}
\newcommand{\ycov}{\pCap_{yy}}
\newcommand{\ccov}{\pCap_{xy}}
%\newcommand{\xcov}{\mat{A}}
%\newcommand{\ycov}{\mat{D}}
%\newcommand{\ccov}{\mat{C}}
\newcommand{\xinfo}{\infomat_{xx}}
\newcommand{\yinfo}{\infomat_{yy}}
\newcommand{\cinfo}{\infomat_{xy}}

\newcommand{\sv}{\statecont}
\newcommand{\pos}{\bm{\rho}}
\newcommand{\vel}{\mathbf{v}}
\newcommand{\acc}{\mathbf{a}}


% 2019 pdf commands
%\newcommand{\pdf}{p}
%\newcommand{\cdf}{P}
%\newcommand{\pmf}{\pdf}

% 2020 pdf commands
\newcommand{\pdf}{p}
\newcommand{\cdf}{P}
\newcommand{\pmf}{p}
\newcommand{\mixture}{f}


\newcommand{\prob}{\mathrm{Pr}}

\newcommand{\transmat}{\mat{F}}
\newcommand{\qdisc}{\mat{Q}}
%\newcommand{\disco}{\Delta} % Discretization time interval t_k - t_{k-1}
\newcommand{\disco}{T} % Discretization time interval t_k - t_{k-1}

\newcommand{\iid}{i.i.d.}

\newcommand{\binopdf}[3]{\binom{#2}{#1}#3^{#1} (1-#3)^{#2-#1}}

\newcommand{\clutterRate}{\Lambda}
\newcommand{\clutterin}{\lambda}
\newcommand{\valgate}{\mathcal{G}}
\newcommand{\assocPDA}[2]{\theta_{#1}^{#2}}
\newcommand{\betaPDA}[2]{\beta_{#1}^{#2}}
\newcommand{\ePDA}[2]{e_{#1}^{#2}}
\newcommand{\lPDA}[2]{l_{#1}^{#2}} % PDAF likelihood ratios in likelihood-ratio formulation of PDAF based on wake-article.
\newcommand{\innovPDA}[2]{\innov_{#1}^{#2}}
\newcommand{\meaPDA}[2]{\mea_{#1}^{#2}}
\newcommand{\xHatPDA}[2]{\hat{\sv}_{#1}^{#2}}
\newcommand{\pHatPDA}[2]{\mathbf{\hat{P}}_{#1}^{#2}}
\newcommand{\kalmanPDA}[2]{\kalmangain_{#1}^{#2}}

\newcommand{\rvx}{X}
\newcommand{\rvy}{Y}

\newcommand{\oSpace}{\mathbb{O}} % Outcome space
\newcommand{\pMeasure}{P} % Probability measure


% IMM commands

\newcommand{\fMM}[1]{\mat{f}^{(#1)}}
\newcommand{\qMM}[1]{\mat{Q}^{(#1)}}
\newcommand{\hMM}[1]{\mat{h}^{(#1)}}
\newcommand{\rMM}[1]{\mat{R}^{(#1)}}
\newcommand{\HMM}[1]{\mathbf{H}^{(#1)}}

\newcommand{\nowMode}{s_k}
\newcommand{\prevMode}{s_{k-1}}
\newcommand{\oldModes}{a}

% Stochastic processes commands

\newcommand{\spr}{\mat{x}}
\newcommand{\spry}{\mat{y}}
\newcommand{\acf}{R}
\newcommand{\sysmat}{\mat{A}}
\newcommand{\pnoisemat}{\mat{G}}
\newcommand{\heaviside}{u}
\newcommand{\wiener}{b}
\newcommand{\psd}{S}
\newcommand{\pnoisecont}{\mat{n}}
\newcommand{\inpucont}{\mat{u}}
\newcommand{\inpu}{\mat{u}}
\newcommand{\inpuscalar}{u}
\newcommand{\inpumat}{\mat{B}}

% IPDA commands

%\newcommand{\exi}{\epsilon}
\newcommand{\exi}{r}
\newcommand{\visi}{\mathrm{v}}
\newcommand{\pobs}{\eta}


% JPDA commands

\newcommand{\tConfig}{\bm{\tau}}
\newcommand{\detEvent}{\tConfig}
\newcommand{\reward}{\mat{A}}
\newcommand{\murtyAssProb}{\mathsf{A}}
\newcommand{\murtyDynList}{\mathsf{L}}
\newcommand{\murtyReturnList}{\mathsf{R}}
\newcommand{\murtyQ}{\mathsf{Q}}
\newcommand{\murtyM}{\mathsf{M}}


% MHT chapter commands

\newcommand{\reidNumber}{\gamma}
\newcommand{\reidConfig}{\tau}
\newcommand{\newtarcount}{\beta}
\newcommand{\ntgt}{N_{\mathrm{TGT}}}

\newcommand{\xpos}{x}
\newcommand{\ypos}{y}
\newcommand{\xvel}{v_x}
\newcommand{\yvel}{v_y}

\newcommand{\class}{C}  % The discrete random estimatee
\newcommand{\classization}{c}  % The discrete random estimatee
\newcommand{\data}{\Delta}  % The generic data
%\newcommand{\data}{Z^k}  % The generic data
%\newcommand{\datazation}{D}  % The generic data
\newcommand{\datazation}[1]{Z^{#1}}  % The generic data
\newcommand{\hindex}{t}  % The generic data
\newcommand{\tindex}{t}  % Target number

\newcommand{\birthVal}{b}
\newcommand{\lambdauVal}{\lambdau}
\newcommand{\indicator}{1}

\newcommand{\assignment}{\psi}
\newcommand{\sregion}{\Omega}
\newcommand{\descent}{d}
\newcommand{\trackscore}{s}
\newcommand{\ntracks}{n}
\newcommand{\mtotal}{M} % Total number of measurements hitherto observed
\newcommand{\linprogsol}{\bm{\tau}}
\newcommand{\uLagrange}{u}
\newcommand{\cost}{c}
\newcommand{\subgradient}{\mat{g}}
\newcommand{\prices}{\mat{p}}

% RFS commands

\newcommand{\sdmarkov}{f_{\dSetization}}
\newcommand{\sdlikeli}{f_{\mSetization}}
\newcommand{\sdpost}[1]{f_{#1}}
\newcommand{\sdpri}[1]{f_{#1}}
\newcommand{\targety}{\statey}
%\newcommand{\measu}{z}
\newcommand{\measu}{\mea}
\newcommand{\wphd}{u} % PHD GM weights, e.g., in \lambdau
\newcommand{\bphd}{b^*} % scale factor for birth-proportional part of  \lambdau

%\newcommand{\phd}{v}
\newcommand{\noexist}{\text{\o}}

% EKF chapter commands

\newcommand{\turnrate}{\omega}
\newcommand{\proposal}{q}
\newcommand{\isfun}{\mat{f}} % Function to be evaluated in importance sampling
\newcommand{\istar}{\pi} % Target density in importance sampling

% ESKF chapter commands

\newcommand{\skewsym}{\bm{S}}
\newcommand{\quat}{\bm{q}}
\newcommand{\quatvec}{\bm{\epsilon}}
\newcommand{\quatscalar}{\eta}
\newcommand{\pquat}{\bm{p}}
\newcommand{\angvel}{\bm{\omega}}
\newcommand{\rotmataa}[2]{\rotmat_{{#1},{#2}}} % Argument 1 is axis. Argument 2 is angle
\newcommand{\gvec}{\mat{g}} % Gravity vector
\newcommand{\errangle}{\bm{\theta}}
\newcommand{\vectoraa}{\mat{v}} % Vector to be rotated in axis angle representation
\newcommand{\nvector}{\mat{n}} % Vector to rotate around in axis angle representation
\newcommand{\nscalar}{n} % Should match \nvector
\newcommand{\angleaa}{\alpha} % Rotation angle in axis angle representation
\newcommand{\rodtrans}{\mat{T}} % Transformation matrix used in proof of Rodrigues formula
\newcommand{\vectorw}{\mat{w}} % Vector to be rotated in axis angle representation
\newcommand{\mvector}{\mat{m}} % Vector to rotate around in axis angle representation
\newcommand{\body}{b} % 
\newcommand{\world}{w} % 
\newcommand{\meacovins}{\mat{R}} % 
\newcommand{\ex}{\mathbf{e}_x}
\newcommand{\SO}{\mathrm{SO}}

% SLAM chapter commands

\newcommand{\map}{\mat{m}}
\newcommand{\odo}{\mat{u}}
\newcommand{\compadd}{\oplus}
\newcommand{\pnoisedisc}{\bm{\Gamma}}
\newcommand{\odox}{u}
\newcommand{\odoy}{v}
\newcommand{\odoyaw}{\varphi}
\newcommand{\idex}{\mat{J}}
\newcommand{\postolm}{\mat{n}}
\newcommand{\delrot}{\mathsf{J}}
\newcommand{\nth}{\will} % Neira-Tardos type hypothesis
\newcommand{\hessian}{\mathsf{H}}
\newcommand{\jacobian}{\mathsf{D}}
\newcommand{\lindex}{i} % SLAM landmark index
\newcommand{\pindex}{\mathsf{p}} % FastSLAM particle index
\newcommand{\slamsv}{\bm{\eta}}
\newcommand{\fxs}{\mat{f}_{\slamsv}}
\newcommand{\xHatSlam}[1]{\bm{\hat{\eta}}_{#1}}
\newcommand{\xBarSlam}[1]{\bm{\hat{\eta}}_{#1|#1-1}}

% Graph-SLAM commands

\newcommand{\winlen}{L}
\newcommand{\slamcost}{s}
\newcommand{\slamerr}{\mat{e}}
\newcommand{\slamjac}{\mat{J}}
%\newcommand{\slamcomplete}{\slamsv_{1:L} }  % Kaess uses \bm{\theta} for this
\newcommand{\slamcomplete}{\bm{\theta} }  % Kaess uses \bm{\theta} for this
\newcommand{\slaminfo}{\bm{\Omega}}
\newcommand{\qmat}{\bm{Q}}  % Q-matrix in the QR decomposition
\newcommand{\rmat}{\bm{R}}  % Q-matrix in the QR decomposition
\newcommand{\newblockrow}{\mat{w}^{\transpose}}  % New block row for incremental ISAM update
\newcommand{\newblockerr}{\gamma}  % New error vector component for incremental ISAM update


% EXAM 2019


\newcommand{\audi}{a}
\newcommand{\pMatOwn}{\mat{U}}
\newcommand{\pMatAudi}{\mat{V}}
\newcommand{\velscalar}{\mathrm{v}}

% Ernesto Macros

%Gaussian
\newcommand{\gaussianBell}[3]{1/(#3*sqrt(2*pi))*exp(-((#1-#2)^2)/(2*#3^2))}
\newcommand{\bell}[4]{#4*exp(-((#1-#2)^2)/(2*#3^2))}
\newcommand{\bellFactorFirstDer}[3]{
    -(#1-#2)/#3^2
}
\newcommand{\bellFactorSecondDer}[3]{
    ((#1-#2)^2-#3^2)/#3^4
}
\newcommand{\normalCDF}[3]{
    (1+erf((#1-#2)/(#3*sqrt(2))))/2
}

%Sampling
\newcommand{\transition}[3]{(-#2*(exp(-#3*#1)-1) +  #3*(exp(-#2*#1)-1))/(#2-#3)}

%Poisson
\newcommand{\poissonPDF}[2]{((#2^#1)*exp(-#2))/(factorial(#1))}
\pgfmathdeclarefunction{poissonCDF}{2}{
    \pgfmathparse{#1>0 ? 0 : 1}
    \ifnum\pgfmathresult=1\pgfmathparse{\poissonPDF{#1}{#2}+poissonCDF(#1-1,#2)}\else\pgfmathparse{0}\fi%
}

% Observability chapter commands

\newcommand{\qobs}{\mat{Q}_O}

% PGM lecture

\newcommand{\facto}{\psi}
\newcommand{\neh}{\mathrm{ne}}
\newcommand{\hmmstate}{s}
\newcommand{\hmmparam}{\theta}
\newcommand{\hmmdata}{z}

% PGM Chapter

\newcommand{\bayesq}{q}  % conditional densities in Bayes net constructed from factor grahp

% Sub-parts are labelled (c1), (c2), ... to match the code, the tests and pytest -m.
\renewcommand{\thesubpart}{\thepartno\arabic{subpart}}
\renewcommand{\subpartlabel}{(\thesubpart)}

% ga2_graded2_student.tex defines \studentversion and builds this file without the solutions.
\ifdefined\studentversion\else\printanswers\fi
\assignmentnr{2}
\deadline{Tuesday {3}. November 23:59}

\begin{document}
{\bf Hand in:} \emph{Tuesday {3}. November 23:59} on Canvas.

\textbf{Groups size:} This assignment can be solved in groups of two or three. You are allowed to complete the assignment on your own if you want, but we advice you to solve it in groups as we believe that group discussions are important. Via Canvas you should also be able to see if there are groups with open spots that you could contact. It is okay to discuss the assignment with other groups, but each group should write an individual report.

\textbf{Group signup:} You find the group signup under \texttt{People} $\rightarrow$ \texttt{Groups} $\rightarrow$ \texttt{GA2 group} in the menu on the left on Canvas. You and your group should sign up in the same group.

\textbf{Hand in:} Hand in the code as \textbf{the zip created by running \textit{create\_handin.py}}. Do not use any libraries that are not part of the TTK4250 environment.
Hand in the report as a single pdf.
Only one student per group should hand in both the report and the code.

\section*{Code content:}
The assignment is pass/fail. The number of tests your code passes is part of how we evaluate the submission as a whole, together with the report.
If you notice the tests having false negatives or positives, or behaving weirdly, please let us know.

\section*{Report content:}
For the report you are limited to 1300 words and a total of 7 pages. Figure captions do not count toward the limit, as long as they are concise and purely descriptive. The group number and group members should be clearly stated.

In the report we require you to clearly answer Task \ref{q:task02}, \ref{q:task03} and \ref{q:task04}. As the report is short you can skip the introduction and conclusion.
You can assume the person reading the report has good knowledge of the subject.
You can simply refer to equations and algorithms in the book when what you want to say is already stated there.
Ex: "Given what's stated in (12.34) and equation (56.78) in the book the oscillations in Figure 9 are expected" (When refering to equations, figures etc. you should make it clear if you are refering to equations, figures etc. in the book or report).

The goal of the report is to show that you understand the code and the theory behind it by presenting and discussing the plots and results. Remember to argue for your claims and explain your reasoning: a statement that is right but not explained does not show understanding.
Ex: If you write "These parameters gave bad results", you should also include something like "as can be seen by the high RMSE values in Figure 9".

It is not expected that you present all the results, plots, and tuning variables from the code, but rather that you present and discuss the ones you find most relevant for each part of the assignment. The report should, however, include all four main types of plots (state, NIS, NEES, error).

In Tasks \ref{q:task02} and \ref{q:task03} you will experiment with different sets of tuning values. Every run saves the figures below in its run directory (re-plot a finished run with \pythoninline{plot_run}), and prints a summary with the RMSE, the averaged NEES and NIS with their 95\% intervals and, on the simulated data, the map-quality counts. You do not need every figure for every set: use the ones that support your discussion, and combine them into multi-panel figures where that saves space.

From the simulated data:
\begin{itemize}
	\item \pythoninline{final_snapshot}: the trajectory and the map against the ground truth, with the map-quality counts in the title.
	\item \pythoninline{pose_diagnostics}: the \(x\), \(y\) and heading errors over time with their \(\pm 1\), 2 and \(3\sigma\) covariance envelopes, and their RMSE.
	\item \pythoninline{pose_nees}: the pose NEES over time with its 95\% bounds and the share inside them, and the ANEES with its interval.
	\item \pythoninline{landmark_nis}: the landmark NIS per degree of freedom over time with its 95\% bounds and the share inside them, and the ANIS with its interval.
\end{itemize}

From Victoria Park, which has no ground truth:
\begin{itemize}
	\item \pythoninline{final_snapshot}: the trajectory and the map against the GPS.
	\item \pythoninline{position_nis}: the NIS of the GPS positions over time with its 95\% bounds and the share inside them, and the ANIS with its interval.
	\item \pythoninline{landmark_nis}: as for the simulated data.
	\item \pythoninline{landmark_growth}: the number of landmarks over time, which shows when data association starts to fail.
\end{itemize}
The timing figures (\pythoninline{timing_*} and \pythoninline{landmarks_and_timing}) are there if run time helps your discussion.

We expect you to discuss what these plots and numbers mean in practice, and how they are connected to the different tuning parameters.

Answers and analysis should connect theory and results to the real world, and show your understanding of the problem and the solution. Theoretical or practical insight from the lectures or book elaborated upon or applied to the datasets counts positive.

Other plots, numbers, and considerations are welcome too, but what to show is \emph{up to you} to decide.
Try to connect the results on the simulated data to the results of the real data where it is possible.

If you have any questions please make a post on the GA2 discussion thread on Canvas and we will try to answer them as soon as possible or contact the student assistants during the excise class.

\noindent\makebox[\linewidth]{\rule{\textwidth}{0.4pt}}

\section*{Some background on the template code}
The template code is a trimmed-down version of a factor graph SLAM system built on \href{https://gtsam.org}{GTSAM}, the library referred to in Sec.\ 9.5 of the book. GTSAM provides the back-end: it holds the factor graph, maintains the Bayes tree, solves for the MAP estimate incrementally with iSAM2, and recovers marginal covariances from it. Everything around that is the front-end: turning odometry and lidar data into factors, predicting the measurements and their innovation covariance, deciding which measurement belongs to which landmark, and when a new landmark is born. Data association itself (JCBB) is given, the factors and the quantities it relies on are yours to implement in Task~\ref{q:task01}.

Set up the environment from the handout folder with
\begin{center}
    \pythoninline{conda env create -f environment.yml} \\
    \pythoninline{conda activate ttk4250_ga2} \\
    \pythoninline{pip install -e . --no-deps}
\end{center}
Run the tests with \pythoninline{pytest}. Runs are started with \pythoninline{run_sim} and \pythoninline{run_real}, and an existing run can be re-plotted with \pythoninline{plot_run <run directory>}. Ctrl+C stops a run and still saves it up to the last completed step; add \pythoninline{--no-save-on-abort} if you would rather it just exit. Setup problems, especially on Windows, are covered in \pythoninline{docs/INSTALL.md}.

The files handed out are:
\begin{itemize}
	\item \pythoninline{configs/}: the run configurations. Everything marked \pythoninline{# TODO tune} is yours to experiment with in Tasks~\ref{q:task02} and~\ref{q:task03}.
	\item \pythoninline{docs/GTSAM.md}: every GTSAM call Task~\ref{q:task01} needs, with what each one means. Read it before you start.
	\item \pythoninline{preprocessing.py}: the odometry and lidar front-end. \pythoninline{relative_pose} and \pythoninline{preintegrate} are yours; the lidar part at the bottom of the file, which turns a raw scan into range-bearing measurements, is given.
	\item \pythoninline{factor_graph.py}: pose prediction, the odometry and landmark factors, the measurement model and its inverse, the joint covariance assembled from GTSAM's marginals, and the innovation covariance. Most of your work is here.
	\item \pythoninline{landmark_manager.py}: landmark birth. The \(M\) of \(N\) rule is yours; the rest is given.
	\item \pythoninline{data_association.py}: the JCBB algorithm, and the forced wrong association for Task~\ref{q:task04}. \textbf{NB:} you do not need to do anything in this file.
	\item \pythoninline{evaluation.py}: the NEES, the NIS, their bounds and the run summary, and, on the simulated data, the comparison of the finished map with the true landmarks. \textbf{NB:} you do not need to do anything in this file.
	\item \pythoninline{slam.py}: the main loop that ties all of the above together. Reading it is the fastest way to understand what happens in one time step. \textbf{NB:} you do not need to do anything in this file.
	\item \pythoninline{loaders/}, \pythoninline{logger.py}, \pythoninline{plotter.py}, \pythoninline{plotting/}: dataset adapters, run logging, and all the plotting. \textbf{NB:} you do not need to do anything in these files.
\end{itemize}

\rule{\textwidth}{0.4pt}

\begin{questions}
	\breaknoans \breakans
	\titledquestion{Implement factor graph SLAM}
\label{q:task01}
A factor graph SLAM system turns sensor data into a graph and solves it. The back-end (iSAM2, through GTSAM) does the solving: it optimizes the graph and recovers covariances from it. Everything around it is yours to implement in this task: turning wheel odometry into relative poses (a, b), adding the factors (c), predicting each measurement and how uncertain that prediction is (d--f), which is what the given JCBB needs to decide which measurement belongs to which landmark, and placing and confirming new landmarks (g, h).

All functions are marked with \pythoninline{# TODO} in the code. The GTSAM calls you need are listed in \pythoninline{docs/GTSAM.md}, and your editor shows their signatures.

\paragraph{Testing your code.}
Each part has its own tests, selected by a pytest marker named after the part: \pythoninline{pytest -m b} runs the tests for (b), and \pythoninline{pytest -m c2} those for (c2). Add \pythoninline{-x} to stop at the first failure. Use them to check each part as you go. \pythoninline{pytest} runs everything, including the end-to-end tests, which need all parts.

\begin{parts}

    % ------------------------------------------------------------------ (a)
    \part Implement \pythoninline{preprocessing.relative_pose}\unskip
, which turns one wheel-encoder reading into a relative pose increment.

    The Victoria Park vehicle reports a forward speed \(v_e\) measured at the left rear wheel and a steering angle \(\alpha\). The encoder sits a distance \(H\) off the centre line, so during a turn it travels along a different arc than the centre of the rear axle. With \(L\) the wheelbase, the distance between the axles, the kinematic bicycle model gives the body twist
    \begin{align*}
        v_b    & = \frac{v_e}{1 - \frac{H}{L}\tan\alpha}, &
        \omega & = \frac{v_b}{L}\tan\alpha,               &
        \xi    & = \begin{bmatrix} v_b & 0 & \omega\end{bmatrix}^T.
    \end{align*}
    Held constant over the time \(\Delta t\) between two odometry readings, this twist integrates \emph{exactly} to \(\Delta T = \mathrm{Exp}(\xi\,\Delta t)\), which follows a circular arc.

    \begin{longhint}
        The exponential map is covered in Sec.\ 6.2.3. \pythoninline{gtsam.Pose2.Expmap} implements it, taking \(\xi\Delta t\) as a numpy array. The vehicle geometry is the \pythoninline{car} argument: \(H\) is \pythoninline{car.H} and \(L\) is \pythoninline{car.L}.
    \end{longhint}
    \begin{solution} \begin{pythoncode}
def relative_pose(
    vel_encoder: float, steer: float, dt: float, car: Car = CAR
) -> gtsam.Pose2:
    tangent_steer = np.tan(steer)
    velocity_body = vel_encoder / (1.0 - (car.H / car.L) * tangent_steer)
    yaw_rate = (velocity_body / car.L) * tangent_steer

    twist_body = np.array([velocity_body, 0.0, yaw_rate], dtype=float)

    return gtsam.Pose2.Expmap(twist_body * dt)
\end{pythoncode}
 \end{solution}

    % ------------------------------------------------------------------ (b)
    \breakans
    \part Implement \pythoninline{preprocessing.preintegrate}\unskip
, which compounds a sequence of relative poses into a single one, with corresponding covariance.

    Wheel odometry arrives far faster than lidar scans. Rather than putting a pose in the graph for every encoder reading, all increments between two scans are compounded into one relative pose and one covariance, which is added to the graph as a single \pythoninline{BetweenFactorPose2}. This is called \emph{preintegration} (Sec.\ 9.1).

    The mean is repeated composition. For the covariance, propagate it through the Jacobians of each composition:
    \[
        \Sigma \leftarrow H_1 \Sigma H_1^T + H_2 \Sigma_i H_2^T,
    \]
    where \(\Sigma_i\) is the covariance of increment \(i\), and \(H_1\) and \(H_2\) are the Jacobians of \pythoninline{compose} with respect to its first and second argument.

    \begin{longhint}
        Start the recursion from the identity pose \pythoninline{gtsam.Pose2()} and a zero \(3\times3\) covariance. \pythoninline{compose} fills in both Jacobians for you if you pass it two arrays preallocated as \pythoninline{np.zeros((3, 3), order="F")}.
    \end{longhint}
    \begin{solution} \begin{pythoncode}
def preintegrate(
    poses: list[gtsam.Pose2],
    covariances: list[np.ndarray],
) -> tuple[gtsam.Pose2, np.ndarray]:
    compounded = gtsam.Pose2.Identity()
    compounded_cov = np.zeros((3, 3))

    for increment, increment_cov in zip(poses, covariances):
        H1 = np.zeros((3, 3), order="F")
        H2 = np.zeros((3, 3), order="F")

        compounded = compounded.compose(increment, H1, H2)
        compounded_cov = H1 @ compounded_cov @ H1.T + H2 @ increment_cov @ H2.T

    return compounded, compounded_cov
\end{pythoncode}
 \end{solution}

    % ------------------------------------------------------------------ (c)
    \breakans
    \part Build the graph. Each factor you add is one term of the negative log posterior (9.4).
    \begin{subparts}
        \subpart Implement \pythoninline{factor_graph.predict_pose}\unskip
, which returns the predicted pose
        \[
            \hat{x}_k = \hat{x}_{k-1} \oplus u_k,
        \]
        the previous pose estimate compounded with the odometry increment \(u_k\), which is expressed in the body frame of \(x_{k-1}\). For a \pythoninline{gtsam.Pose2}, \(\oplus\) is \pythoninline{compose}.

        Unlike the prediction step of an EKF, this prediction does not alter the estimate. It only gives the optimizer a starting guess for the new pose \(x_k\), which an iterative solver needs for every variable. What connects \(x_k\) to \(x_{k-1}\) is the odometry factor you add in (c2).
        \begin{solution} \begin{pythoncode}
def predict_pose(previous_pose: gtsam.Pose2, relative_pose: gtsam.Pose2) -> gtsam.Pose2:
    return previous_pose.compose(relative_pose)
\end{pythoncode}
 \end{solution}

        \subpart Implement \pythoninline{factor_graph.add_odometry_factor}\unskip
, adding the middle term of (9.6) as a \pythoninline{BetweenFactorPose2}.
        \begin{longhint}
            The factor is constructed as
            \begin{center}\pythoninline{gtsam.BetweenFactorPose2(key_previous, key_current, measured, noise_model)}\end{center}
            and \pythoninline{graph.add(factor)} adds it to the graph. Use \pythoninline{gtsam.noiseModel.Gaussian.Covariance}, which takes the full \(3\times3\) matrix, and not \pythoninline{.Sigmas}, which takes only standard deviations. The covariance you preintegrated in (b) is not diagonal. A heading error also moves the pose sideways, so the two are correlated. Dropping that correlation makes the factor too confident in some directions and not confident enough in others.
        \end{longhint}
        \begin{solution} \begin{pythoncode}
def add_odometry_factor(
    graph: gtsam.NonlinearFactorGraph,
    key_previous: int,
    key_current: int,
    relative_pose: gtsam.Pose2,
    relative_pose_cov: np.ndarray,
) -> None:
    noise = gtsam.noiseModel.Gaussian.Covariance(relative_pose_cov)
    graph.add(
        gtsam.BetweenFactorPose2(
            key_previous,
            key_current,
            relative_pose,
            noise,
        )
    )
\end{pythoncode}
 \end{solution}

        \subpart Implement \pythoninline{factor_graph.add_landmark_factor}\unskip
, adding the last term of (9.6) as a \pythoninline{BearingRangeFactor2D}.
        \begin{longhint}
            The factor is constructed as
            \begin{center}\pythoninline{gtsam.BearingRangeFactor2D(pose_key, landmark_key, bearing, range, noise_model)}\end{center}
            with the bearing \emph{first}, as a \pythoninline{gtsam.Rot2(angle)}, and the range as a float. A plain float bearing raises a \pythoninline{TypeError}. The noise model you are handed is already ordered [bearing, range] to match.
        \end{longhint}
        \begin{solution} \begin{pythoncode}
def add_landmark_factor(
    graph: gtsam.NonlinearFactorGraph,
    pose_key: int,
    landmark_key: int,
    measurement: np.ndarray,
    noise_model: gtsam.noiseModel.Base,
) -> None:
    measured_range, measured_bearing = float(measurement[0]), float(measurement[1])
    graph.add(
        gtsam.BearingRangeFactor2D(
            pose_key,
            landmark_key,
            gtsam.Rot2(measured_bearing),
            measured_range,
            noise_model,
        )
    )
\end{pythoncode}
 \end{solution}
    \end{subparts}

    % ------------------------------------------------------------------ (d)
    \breakans
    \part Implement \pythoninline{factor_graph.predict_measurement}\unskip
, the measurement function \(\mathbf{h}(x_k, m^i)\) of (9.6), together with the two blocks of the Jacobian \(\mathbf{J}\) in (9.8) belonging to this pose and this landmark. It returns the predicted measurement \(z = [\text{range}, \text{bearing}]\), the pose block \(\partial z / \partial x_k\) (\(2\times3\)) and the landmark block \(\partial z / \partial m^i\) (\(2\times2\)). You do not derive any of this yourself. GTSAM computes the values and the Jacobians, one row at a time. Your task is to assemble them:
    \begin{enumerate}
        \item Preallocate four sub-Jacobians, one row each: range and bearing, each with respect to the pose (\(1\times3\)) and to the landmark (\(1\times2\)), e.g.\ \pythoninline{H_range_pose = np.zeros((1, 3), order="F")}. Any other shape or type makes GTSAM raise a \pythoninline{TypeError}.
        \item Pass them to the two GTSAM calls. Each returns its value and fills in its two sub-Jacobians:
        \begin{center}
            \pythoninline{pose.range(landmark, H_range_pose, H_range_landmark)} \\
            \pythoninline{pose.bearing(landmark, H_bearing_pose, H_bearing_landmark)}
        \end{center}
        \pythoninline{bearing} returns a \pythoninline{gtsam.Rot2}, not a float; take \pythoninline{.theta()} for the angle.
        \item Stack the rows with \pythoninline{np.vstack}, range first and bearing second, to match the order of \(z\).
    \end{enumerate}
    \begin{solution} \begin{pythoncode}
def predict_measurement(
    pose: gtsam.Pose2,
    landmark: np.ndarray,
) -> tuple[np.ndarray, np.ndarray, np.ndarray]:
    landmark = gtsam.Point2(float(landmark[0]), float(landmark[1]))

    H_range_pose = np.zeros((1, 3), order="F")
    H_range_landmark = np.zeros((1, 2), order="F")
    H_bearing_pose = np.zeros((1, 3), order="F")
    H_bearing_landmark = np.zeros((1, 2), order="F")

    predicted_range = pose.range(landmark, H_range_pose, H_range_landmark)
    predicted_bearing = pose.bearing(
        landmark, H_bearing_pose, H_bearing_landmark
    ).theta()

    z = np.array([predicted_range, predicted_bearing])
    H_pose = np.vstack((H_range_pose, H_bearing_pose))
    H_landmark = np.vstack((H_range_landmark, H_bearing_landmark))

    return z, H_pose, H_landmark
\end{pythoncode}
 \end{solution}

    % ------------------------------------------------------------------ (e)
    \breakans
    \part Implement \pythoninline{factor_graph.assemble_joint_covariance}\unskip
, which builds the joint covariance \(P\) of the current pose and the landmarks in the local map: the \(P\) in \(S = HPH^T + R\) of (f).

    A factor graph stores the sparse information matrix \(\Lambda\), not the covariance. GTSAM computes \(P\) from \(\Lambda^{-1}\) for the variables you ask for (Sec.\ 9.4). It hands you \(P\) as a \pythoninline{gtsam.JointMarginal}, in blocks:
    \begin{center}\pythoninline{joint_marginal.at(key_i, key_j)}\end{center}
    is the covariance between the two variables \pythoninline{key_i} and \pythoninline{key_j}. Your task is to put these blocks together into one matrix, in the order of \pythoninline{keys}: the pose first, then the landmarks. Either stack the blocks with \pythoninline{np.block}, or preallocate \(P\) and fill it block by block in two loops over \pythoninline{keys}. For the loops, the size of each variable's block is \pythoninline{joint_marginal.at(key, key).shape[0]}.

    Be sure to include every block, also the ones pairing two different landmarks, \(P_{m^i m^j}\). Leaving them out gives a plausible-looking but wrong \(S\). Also keep the order of \pythoninline{keys}, since (f) stacks the Jacobians in that order.

    \begin{solution} \begin{pythoncode}
def assemble_joint_covariance(
    joint_marginal: gtsam.JointMarginal,
    keys: list[int],
) -> np.ndarray:
    return np.block([
        [joint_marginal.at(key_i, key_j) for key_j in keys]
        for key_i in keys
    ])

    # Equivalent, with two loops filling a preallocated P:
    #
    # dims = [joint_marginal.at(key, key).shape[0] for key in keys]
    # offsets = np.concatenate([[0], np.cumsum(dims)])  # where each block starts
    # P = np.zeros((offsets[-1], offsets[-1]))
    # for i, key_i in enumerate(keys):
    #     for j, key_j in enumerate(keys):
    #         rows = slice(offsets[i], offsets[i + 1])
    #         cols = slice(offsets[j], offsets[j + 1])
    #         P[rows, cols] = joint_marginal.at(key_i, key_j)
    # return P
\end{pythoncode}
 \end{solution}

    % ------------------------------------------------------------------ (f)
    \breakans
    \part Implement \pythoninline{factor_graph.innovation_covariance}\unskip
, which assembles
    \[
        S = H P H^T + R
    \]
    for the whole local map at once. This is (9.26) with \(P\) the joint covariance assembled in part (e), \(H\) the stacked Jacobian \(\mathbf{J}_a\) from part (d), and \(R\) the block-diagonal measurement noise \(\mathbf{R}_a\).

    Here \(n\) = \pythoninline{len(jacobians_pose)} is the number of landmarks in the local map. Keep the dimensions in mind: \(P\) is \((3+2n)\times(3+2n)\), pose first and then each landmark; \(H\) is \(2n\times(3+2n)\), two rows per predicted measurement; \(R\) and the returned \(S\) are \(2n\times2n\). For \(n = 3\):
    \[
        H = \bordermatrix{
                & x_k   & m_0   & m_1   & m_2   \cr
            z_0 & H^x_0 & H^m_0 & 0     & 0     \cr
            z_1 & H^x_1 & 0     & H^m_1 & 0     \cr
            z_2 & H^x_2 & 0     & 0     & H^m_2 \cr
        }
    \]
    where \(H^x_i\) is \pythoninline{jacobians_pose[i]} (\(2\times3\)) and \(H^m_i\) is \pythoninline{jacobians_landmark[i]} (\(2\times2\)). \(H\) is dense in its three leftmost columns, because every predicted measurement depends on the pose, and block diagonal elsewhere, because predicted measurement \(i\) depends only on landmark \(i\). \(R\) is \pythoninline{measurement_cov} repeated \(n\) times along its diagonal, since the measurements are independent.

    \begin{solution} \begin{pythoncode}
def innovation_covariance(
    jacobians_pose: list[np.ndarray],
    jacobians_landmark: list[np.ndarray],
    joint_covariance: np.ndarray,
    measurement_cov: np.ndarray,
) -> np.ndarray:
    n = len(jacobians_pose)

    H = np.zeros((2 * n, 3 + 2 * n))
    for i in range(n):
        H[2 * i : 2 * i + 2, 0:3] = jacobians_pose[i]
        H[2 * i : 2 * i + 2, 3 + 2 * i : 3 + 2 * i + 2] = jacobians_landmark[i]

    R = np.kron(np.eye(n), measurement_cov)

    return H @ joint_covariance @ H.T + R
\end{pythoncode}
 \end{solution}

    % ------------------------------------------------------------------ (g)
    \breakans
    \part Implement \pythoninline{factor_graph.inverse_measurement}\unskip
, the inverse measurement function \(\mathbf{h}^{-1}(z, x)\): given a measurement \(z\) taken from pose \(x\), it returns the landmark's position in the world frame. It places a new landmark for a measurement that JCBB could not associate. Such a measurement may also be clutter, so (h) decides when the new landmark enters the graph.

    Unlike in EKF-SLAM (Sec.\ 7.2), only the position is needed, not a covariance: once the landmark is connected by factors, its covariance follows from the graph, like that of every other variable.

    \begin{longhint}
        Find the landmark in the body frame first. \pythoninline{rotate} needs a point, not a range, so start from \pythoninline{gtsam.Point2(range, 0.0)}, the landmark straight ahead, and let \pythoninline{gtsam.Rot2(bearing).rotate(...)} turn it to the measured bearing, giving \((r\cos b,\ r\sin b)\). Then \pythoninline{pose.transformFrom(...)} takes that point from the body frame to the world frame.
    \end{longhint}
    \begin{solution} \begin{pythoncode}
def inverse_measurement(pose: gtsam.Pose2, measurement: np.ndarray) -> np.ndarray:
    measured_range, measured_bearing = float(measurement[0]), float(measurement[1])
    landmark_body = gtsam.Rot2(measured_bearing).rotate(
        gtsam.Point2(measured_range, 0.0)
    )
    landmark_world = pose.transformFrom(landmark_body)
    return np.asarray(landmark_world, dtype=float).reshape(2)
\end{pythoncode}
 \end{solution}

    % ------------------------------------------------------------------ (h)
    \part Implement \pythoninline{landmark_manager.TentativeLandmark.is_confirmed}\unskip
, the \(M\) of \(N\) confirmation rule from target tracking: confirm a tentative landmark once it has been seen in at least \(M\) of the last \(N\) time steps, the current one included.

    The evidence is the list \pythoninline{self.supporting_observations}, with one entry per time step the landmark was seen in. The given code fills it, starting with the measurement that created the landmark. Each entry's time step is in \pythoninline{.step}. Your function only reads the list: count the entries whose \pythoninline{.step} lies inside the window, and compare with \(M\).

    \begin{longhint}
        Both ends are easy to get off by one. With \(N = 3\) and \pythoninline{current_step} 10, the window is steps 8, 9 and 10: it starts at \(10 - 3 + 1\), and a step equal to the start counts, so compare with \(\geq\). Exactly \(M\) hits is enough, so compare the count with \(\geq\) as well.
    \end{longhint}
    \begin{solution} \begin{pythoncode}
def is_confirmed(self, current_step: int, M: int, N: int) -> bool:
    window_start = current_step - N + 1
    hits_in_window = sum(
        observation.step >= window_start
        for observation in self.supporting_observations
    )
    return hits_in_window >= M
\end{pythoncode}
 \end{solution}

\end{parts}

	\titledquestion{Run factor graph SLAM on simulated data}
\label{q:task02}
You are given a data set consisting of
\begin{itemize}
    \item odometry (K, 3): the measured movement of the robot in body frame
    \item z (K,): Python list of detections at each time step. Each detection is (\#measurements, 2).
    \item poseGT (K, 3): the pose ground truth at each time step
    \item landmarks (M, 2): the ground truth landmarks
\end{itemize}

We have given you initial tuning values in \pythoninline{configs/sim_default.yaml}; everything marked \pythoninline{# TODO tune} is yours to change. They are a \emph{working baseline}: the run completes, but the estimate is neither accurate nor consistent. The pose NEES sits well above its expected value, and the landmark NIS well below it. With a correct Task~\ref{q:task01}, these values map 72 of the 78 true landmarks, with 3 spurious ones, and give a pose ANEES of 8.37; the run is deterministic, so if yours differs, look at Task~\ref{q:task01} before you start tuning. We want you to find 3 different sets of tuning values that differ in different ways and explain the results you get using them.

Every run prints a summary of the averaged NEES and NIS (the landmark ANIS per degree of freedom, since the number of associated landmarks varies). To judge consistency, compare each average with its expected value, and check that about 95\% of the per-step values fall inside their bounds. Do not rely on the 95\% interval of the average alone: (4.67) assumes independent errors, but the errors of neighbouring poses in a smoother are strongly correlated, so even with the correct noise, one run's ANEES can land anywhere between about 1.5 and 8.

An example of three different sets of tuning values could be
\begin{itemize}
    \item a set that gives better performance than the initial tuning values in some way
    \item a set that gives similar performance in terms of error but bad performance in terms of consistency
    \item a set that maps noticeably fewer landmarks, and how that affects the trajectory
\end{itemize}

\hfill

Some questions to help discussion (not mandatory to answer):
\begin{itemize}
    \item Which types of errors do you see?
    \item Why do you think the errors are like this?
    \item Every simulated run reports a map-quality line. Each map landmark is paired with at most one true landmark, closest pairs first and only within 3\,m, and only the true landmarks the run actually observed are counted:
    \begin{itemize}
        \item \emph{Mapped}: a true landmark with a map landmark paired to it. The RMS position error is taken over these pairs.
        \item \emph{Missed}: an observed true landmark with no map landmark paired to it. Either it was never confirmed, because it was not seen in \(M\) of \(N\) steps or its measurements kept being associated elsewhere, or its estimate drifted more than 3\,m away.
        \item \emph{Duplicate}: an extra map landmark within 3\,m of a true landmark that already has its pair. Data association failed to match a measurement to the landmark already in the map, and a new landmark was born from it instead.
        \item \emph{Spurious}: a map landmark with no true landmark within 3\,m. It was born from a wrong measurement or a wrong association, or its estimate has drifted far from where the landmark is.
    \end{itemize}
    Duplicates point at the front-end (association and landmark birth); missed and spurious landmarks can be either. Use these numbers to judge how much of your error the front-end is responsible for, and how much the back-end. (Sec.\ 9.1)
    \item Would you say that there are any shortcomings of the algorithm on this data set?
\end{itemize}

\textit{Note}: The runtime of JCBB grows quickly with the number of duplicate landmarks, because each one gives a measurement more candidates to pair with. A tuning that creates duplicates (noise values the gates cannot cope with, or landmark birth on a single sighting, \(M = N = 1\), with a wrong noise model) can slow a run from seconds to minutes. Watch the map-quality line and how many new landmarks appear.
The local map, \pythoninline{sensor.range_local} and \pythoninline{sensor.bearing_local_deg}, is not a tuning value: it is set from what the sensor can see, with a margin. Making it smaller hides landmarks the sensor does see from data association, so their measurements become duplicate landmarks, and the run can slow down to the point of hanging.

	\breaknoans
	\titledquestion{Run factor graph SLAM on the Victoria Park dataset}
\label{q:task03}
The Victoria Park data set is a SLAM data set available along with some additional information
\href{http://www-personal.acfr.usyd.edu.au/nebot/victoria_park.htm}{here}
\footnote{\href{http://www-personal.acfr.usyd.edu.au/nebot/victoria_park.htm}{http://www-personal.acfr.usyd.edu.au/nebot/victoria\_park.htm}}.
All times are converted to seconds for you in the loader, and the functions that create odometry and detections from the raw data are provided.

The dataset consists of
\begin{itemize}
    \item LASER (mK, 361): Laser scanner returns. There are 361 returns per scan evenly divided in a half circle.
    \item speed (K,): measured speed from a wheel encoder at the rear left wheel.
    \item steering (K,): measured wheel steering angle
    \item Lo\_m (Kgps,): GPS longitude in meters
    \item La\_m (Kgps,): GPS latitude in meters
    \item timeLsr (mK,): Time of laser scans in milliseconds
    \item timeOdo (K,): Time of odometry measurements in milliseconds
    \item timeGps (Kgps,): Time of GPS measurements in milliseconds
\end{itemize}

As in the previous task, we have given you initial tuning values in \pythoninline{configs/real_default.yaml}. They are a \emph{working baseline}: the map is clean, loops close, and the trajectory stays close to the GPS. They are, however, not consistent: the landmark NIS sits well below its expected value. Finding out why is the main point of this task.

In the same way as previously, we want you to find 3 different sets of tuning values that differ in different ways and explain the results you get using them. At least one of the sets should lower the odometry noise to bring the landmark NIS closer to its expected value. Describe what happens to the map and to the trajectory compared with the GPS, and explain why. It helps to recall what your \pythoninline{preintegrate} from Task~\ref{q:task01} assumes about consecutive odometry increments, and to compare that with how wheel-odometry errors behave on a real car.

\textbf{Plan your runs.} A full run over the data set (the default: all 7247 laser scans) takes a few minutes. It took about 2.5 minutes on our machine with the given tuning, and about twice that with a tuning that makes the map grow, such as lowered odometry noise. Expect it to take longer on a laptop. To find good values without spending hours waiting:
\begin{itemize}
    \item Explore with \pythoninline{run_real --steps 2000}, which takes about 20 seconds. The first 2000 scans cover about a quarter of the drive and most of the park, including the first loop closures, and a change in tuning moves the results in the same direction as it does on a full run.
    \item Change one group of parameters at a time (odometry noise, measurement noise, gates, landmark birth), so that you know what caused what.
    \item Watch the number of landmarks. The given tuning ends the first 2000 scans with about 140. Several hundred means data association is failing, and that tuning will only get worse, and slower, over the full run.
    \item Give every run its own \pythoninline{--output-dir}, and re-plot a finished run with \pythoninline{plot_run} instead of running it again.
\end{itemize}
All figures and numbers in your report should come from full runs: after scan 5000 the car revisits earlier parts of the park, and a short run never gets there.

See the linked webpage and the files in the Victoria Park folder for further information.
The same reasoning applies to NIS and tuning here as in the simulated dataset.

Note that you can use the GPS as a sort of ground truth, but it is very much less than perfect.
You might want to use a NIS calculation instead of NEES calculation if you do that.
Some comparison to the GPS is, however, expected.
The position NIS plot treats each GPS sample as a position measurement with standard deviation \pythoninline{noise.sigma_gnss} (1\,m by default). This value is an assumption about the GPS, not a filter parameter: it never enters the graph, but it sets the scale of that plot, so keep it in mind when you interpret it.

\hfill

Some questions to help discussion (not mandatory to answer):
\begin{itemize}
    \item Which types of errors do you see?
    \item Is this different on the real data set compared to the simulated one? Why do you think this is so?
    \item What happens to the estimate of the whole trajectory when a loop is closed?
    \item Would you say that there are any shortcomings of the algorithm on this data set?
    \item What is your general take on graph-based SLAM on data like this? The book claims in Sec.\ 9.2 that (9.13) is "significantly more robust and accurate than EKF-SLAM" (Sec.\ 7.2). Which of your results speak to that claim, and why?
    \item Can you think of any improvement strategies?
\end{itemize}

	\titledquestion{Two short experiments}
\label{q:task04}
Both of these are quick to run and belong in the report as a figure and a paragraph each. Use whichever data set makes the point most clearly.

\begin{parts}
    \part \textbf{Relinearization.} iSAM2 does not relinearize everything at every step; it relinearizes the variables whose estimates have moved enough to warrant it. Two knobs control this: \pythoninline{backend.relinearize_threshold}, how far an estimate must move, and \pythoninline{backend.relinearize_skip}, how many steps pass between checks.

    Push both to extremes and report what happens to accuracy, consistency and runtime. Relate your results to Sec.\ 9.5.1, which explains why ISAM had to linearize as seldom as possible and why being able to restrict relinearization to the affected variables is what makes ISAM2 worthwhile.

    \part \textbf{One bad association.} Loosening the gates does not make JCBB err on these data sets, so force a mistake instead: \pythoninline{association.inject_wrong_at_step} re-pairs one measurement with the wrong landmark at that scan step (600 works well on the simulated data, 1500 on Victoria Park). Show what this single wrong association does to the map.

    Then set \pythoninline{backend.robust_kernel: huber} and run the same thing again. Explain what the robust kernel does to the cost function (9.6), why that helps, and what it costs you when the associations are in fact all correct.

    In your discussion, compare this with EKF-SLAM (Sec.\ 7.2): what are your options for recovering from a wrong association in each of the two frameworks?
\end{parts}

\hfill

\textit{Optional, if you have time}: \pythoninline{backend.solver: batch} re-solves the entire graph with Levenberg-Marquardt at every step instead of updating it incrementally. It is the honest implementation of (9.13), and iSAM2 is an incremental approximation of it. Run both on the simulated data set and compare the estimates and the runtimes. Do not try it on Victoria Park unless you have an afternoon to spare -- which is itself the point of Sec.\ 9.3.3.

\end{questions}
\end{document}
